\documentclass[12pt]{article}
\usepackage[utf8]{inputenc}
\usepackage[T1]{fontenc}
\usepackage{hyperref}
\title{Boro Language Resources for K-12 Curriculum Development}
\author{Compiled by Claude Sonnet 5.5 (claude-sonnet-5-5) for Andr\'as Kornai}
\date{}
\begin{document}
\maketitle

\section{Introduction}
Boro (ISO 639-3 xxb; Glottocode boro1283) is an extinct language of central-eastern Ghana, formerly spoken around Worawora and Tapa. It is known only from about 12 lexical items collected by Plehn around the end of the 19th century. By then it had reportedly been reduced to a single elderly speaker. Its classification is uncertain: Westermann (1922) placed it among the Togorestsprachen, Glottolog follows him, and Heine (1968) left it unclassified. It should not be confused with the Boro (Bodo) language of India.

\section{Primary Educational Resources}
No primers, textbooks, school programmes or curriculum materials were found in the search results.

\section{Linguistic and Reference Materials}
\begin{itemize}
\item Wikipedia, Boro language (Ghana), which summarises the evidence and classification: \url{https://en.wikipedia.org/wiki/Boro_language_(Ghana)}
\item Wikipedia, Ghana--Togo Mountain languages, the proposed wider grouping: \url{https://en.wikipedia.org/wiki/Ghana%E2%80%93Togo_Mountain_languages}
\end{itemize}

\section{Religious and Community Resources}
None were found. No scripture translations or community organisations for Boro appeared in the search results.

\section{Digital and Audio Resources}
No audio, corpus or software resources were found.

\section{Curriculum Development Context}
With only about 12 recorded words, a curriculum taught in Boro is not possible. A school could instead offer a heritage-awareness unit on the history of the Worawora and Tapa area and on language loss in Ghana. A general article on safeguarding Ghana's indigenous languages is listed below as context.
\begin{itemize}
\item ``Speak It. Preserve It. Pass It On. Why Ghana Must Save Its Indigenous Languages'' (ModernGhana): \url{https://www.modernghana.com/news/1514755/speak-it-preserve-it-pass-it-on-why-ghana-must.html}
\end{itemize}
\end{document}